\documentclass[10pt,a4paper]{amsart}
\usepackage[margin=19mm]{geometry}
\usepackage{amsmath,amssymb,mathtools}
\usepackage[scr=rsfs,cal=euler]{mathalfa}
\usepackage{xfrac}
\usepackage[unicode,hidelinks,bookmarks=false]{hyperref}
\newcommand{\ie}{i.e.\ }
\newcommand{\cf}{cf.\ }
\newcommand{\resp}{respectively\ }
\newcommand{\Subsection}{\subsection}
\providecommand{\nobreakdash}{\nobreak\hbox{-}}
\setlength{\emergencystretch}{2em}
\multlinegap0pt
\usepackage{amssymb}
\newtheorem{prop}{Proposition}
\newtheorem{claim}{Claim}
\newtheorem{lemm}{Lemma}
\title{Universally defined cycles I}
\author{Claire Voisin}
\date{Source: arXiv:2408.06893v2 (CC BY 4.0)}
\begin{document}
\maketitle

\noindent\emph{Source and licence.} Universally defined cycles I. Version arXiv:2408.06893v2 (CC BY 4.0), distributed by arXiv under the Creative Commons Attribution 4.0 International licence, http://creativecommons.org/licenses/by/4.0/, as recorded in the arXiv licence metadata for this exact version and shown on the abstract page https://arxiv.org/abs/2408.06893v2. Full licence notice: \texttt{licenses/arxiv-2408.06893-CC-BY-4.0.txt}.\par\smallskip

\noindent\textit{Editorial scope.} Counted unit: the article's prop lefirstred (source0.tex lines 229--232). The article's complete original proof of it is delivered with it (source0.tex lines 233--288, one \texttt{proof} environment of 56 lines). No mathematics is written, repaired or re-derived: the statement and the proof are the article's own lines, in the article's own notation, and no retained line is substituted or re-flowed.  No further source-local context is needed: every cross-reference of every retained line resolves inside the two delivered spans.  Nothing was omitted from any retained span, so the package carries no omission record.  No retained line contains a citation, so the wrapper carries no bibliography and no bibliography entry.  No retained line contains a figure environment, an image-inclusion command or drawing code, so no image is delivered and the package declares no asset dependency.  Editorial formatting only, in addition: an AMS \texttt{amsart} preamble that restates the article's own declarations and macros used by the retained lines, namely the environments \texttt{prop}, \texttt{claim}, \texttt{lemm} on separate counters, together with the standard packages those lines need.  The retained environments print in this document as Proposition 1, Claim 1, Lemm 1 in order of appearance, whereas the article prints the same items with the numbering of its own full document.  The complete native source of the article as distributed in the arXiv e-print for this exact version is bundled unchanged as \texttt{sources/163729/source0.tex}.
 \begin{prop}\label{lefirstred} If ${\rm dim}\,W$ is large compared to ${\rm dim}\,X$,
 the pull-back and restriction composite map
 ${\rm CH}(X)\stackrel{p_X^*}{\rightarrow }{\rm CH}(G\times X)\rightarrow {\rm CH}((G\times X)\setminus \mathcal{C}_X)$ is injective.
\end{prop}

\begin{proof} Let  $h\in{\rm CH}(\mathbb{P}(E^*))$ be the first Chern class of the line bundle $\mathcal{O}_{\mathbb{P}(E^*)}(1)$.
We observe that
$\mathcal{C}$ is the zero-set of the tautological section of $\mathcal{Q}\boxtimes \mathcal{O}_{\mathbb{P}(E^*)}(1)$ on $G\times \mathbb{P}(E^*)$. Furthermore, as
$W$ is generated by sections, $\mathcal{C}$ is smooth of the right codimension $k$: more precisely,
via the projection to $\mathbb{P}(E^*)$, $\mathcal{C}$ is fibered over $\mathbb{P}(E^*)$ in Grassmannians $G(k,W^{\perp e})$ for hyperplanes $W^{\perp e}\subset W$, hence it is smooth and furthermore the restriction map
${\rm CH}(G\times\mathbb{P}(E^*))\rightarrow {\rm CH}(\mathcal{C})$ is surjective.
We thus conclude that the class $C$ of $\mathcal{C}$ in ${\rm CH}(G\times \mathbb{P}(E^*))$ is
$c_k(\mathcal{Q}\boxtimes \mathcal{O}_{\mathbb{P}(E^*)}(1))$ and that, denoting $j$ the inclusion of
$\mathcal{C}$ in $ G\times \mathbb{P}(E^*)$, ${\rm Im}\,(j_*:{\rm CH}^*(\mathcal{C})\rightarrow
{\rm CH}^{*+k}(G\times \mathbb{P}(E^*)))$ is equal to
${\rm Im}\,(C:{\rm CH}^*(G\times \mathbb{P}(E^*)) \rightarrow
{\rm CH}^{*+k}(G\times \mathbb{P}(E^*)))$.
Finally, as we are working with rational coefficients, the natural map
$n_*:{\rm CH}_*(\mathcal{C})\rightarrow {\rm CH}_*(\mathcal{C}_X)$ is surjective, and thus we conclude, using the localization exact sequence, that
${\rm Ker}\,( {\rm CH}(G\times X)\rightarrow  {\rm CH}((G\times X)\setminus \mathcal{C}_X))$ is equal to
\begin{eqnarray}\label{eqmap} {\rm Im}\,(\pi_*\circ C:{\rm CH}^*(G\times \mathbb{P}(E^*)) \rightarrow
{\rm CH}(G\times X)^{*+1}).
\end{eqnarray}
The space ${\rm CH}^*(G\times \mathbb{P}(E^*))$ is generated over the $\mathbb{Q}$-algebra ${\rm CH}^*(G\times X)$ by the classes $h^j$, for $j=0,\ldots, k-1$. As $\pi_*$ is a
${\rm CH}^*(G\times X)$-linear morphism by the projection formula, we conclude that
$ {\rm Im}\,(\pi_*\circ C)$ is generated over ${\rm CH}^*(G\times X)$ by the classes
$\pi_*(C\cdot h^j)$.  Proposition \ref{lefirstred} thus  follows from the following
\begin{claim}\label{claim} If $\alpha\in{\rm CH}(X)\subset  {\rm CH}(X\times G)$ has the property that $\alpha$
belongs to the ideal generated by the elements  $\pi_*(C\cdot h^j)$, $j=0,\ldots,\,k-1$, then
$\alpha=0$.
\end{claim}
Let us now prove the claim.  We use the same notation  $c_i(\mathcal{Q})$ and $h$ for the pull-back of these classes to
$G\times \mathbb{P}(E^*)$. Then $C=\sum_{i=0}^{i=k} c_i(\mathcal{Q})h^{k-i}$
and it follows that

\begin{eqnarray}\label{eqclasses} \pi_*(C\cdot h^j)=\sum_{i=0}^{i=k} c_i(\mathcal{Q})s_{j-i+1}(E^*)
\end{eqnarray}
for $j=0,\ldots,k-1$, where $s_i(E^*)$ denotes the $i$-th Segre class of $E^*$.
In particular,
\begin{eqnarray}\label{eqclasses5}\pi_*(C\cdot h^j)=c_{j+1}(\mathcal{Q})+P_j,\end{eqnarray} where $P_j$ is a polynomial
in the variables  $c_i(\mathcal{Q})$ with $i\leq j$, with coefficients in ${\rm CH}(X)$, and $P_0\in{\rm CH}(X)$.

 As ${\rm dim}\,W\gg {\rm dim}\,X$ and we can obviously restrict ourselves to  considering cycles of degree
$\leq {\rm dim}\,X$, there are no nontrivial relations in these degrees between the classes
$c_j(\mathcal{Q})$ for $j>0$. In other words, we can do as if ${\rm CH}(G)$ were the
 polynomial ring with generators
$Y_1=c_1(\mathcal{Q}),\ldots,\,Y_k=c_k(\mathcal{Q})$ and thus ${\rm CH}(G\times X)={\rm CH}(X)[Y_1,\ldots, Y_k]$. We know by (\ref{eqclasses5}) that
$\pi_*(C\cdot h^j)=Y_{j+1}+P_j$, where $P_j$ is a polynomial
in the $Y_{i}$ with $i\leq j$.
We apply  then the following elementary lemma to $R={\rm CH}(X)$, $a=\alpha$:
\begin{lemm} Let $R$ be a unitary commutative ring and
let $G_j=Y_j+P_j(Y_1,\ldots,Y_{j-1})\in R[Y_1,\ldots,\,Y_k]$, for $j=1,\ldots, k$. Then
for any relation $\sum_jF_j G_j=a$, with
$F_j\in R[Y_1,\ldots,\,Y_k]$ and  $a\in R$, one has $a=0$.
\end{lemm}
\begin{proof} Indeed, from the structure of the polynomials
$G_j$, we  construct inductively  $r_1,\ldots,r_k\in R$ such that
$G_j(r_1,\ldots, r_k)=0$ for all $j$, by the formulas $r_j=-P_j(r_1,\ldots,r_{j-1})$, $r_0=P_0$. Then $\sum_jF_j(r_\cdot)G_j(r_\cdot)=0=a$.
\end{proof}
Claim \ref{claim} is now proved and the proof of Proposition \ref{lefirstred} is now finished.
\end{proof}

\end{document}
